% ====================================================================
\chapter{Sequences and Limits}\label{ch:limits}
% ====================================================================

\begin{chapterintro}
Economics is full of processes that unfold over time: output in successive quarters, the balance of a savings account
year after year, the price in successive rounds of an auction. A sequence records such a process, and its limit
describes where the process is heading in the long run. This chapter defines sequences, gives the precise definition of
the limit of a sequence --- one of the most important definitions in the whole module --- and shows how to prove that a
sequence converges to a given number, or that it has no limit at all.
\end{chapterintro}

\begin{prerequisites}
Quantifiers and their order (\cref{sec:quantifiers}); proof by contradiction (\cref{sec:proof-types}); functions
(\cref{sec:function}); absolute value as distance and the triangle inequality (\cref{ch:abs}); solving inequalities
(\cref{ch:ineq}).
\end{prerequisites}

\section{Sequences}\label{sec:sequences}

\begin{definition}[Sequence][def:sequence]
A \term{sequence} is an infinite ordered list $x=(x_1,x_2,x_3,\dots)$, whose elements $x_n$ are indexed by $n\in\N$.
More formally, a sequence $x$ is a function whose domain is the set of natural numbers; we write $x_n=x(n)$.
\src{L2, slide 22}
\end{definition}

\begin{example}[Sequences in applications]
\begin{enumerate}
  \item Writing H for heads and T for tails, a sequence of coin-flip results might appear as
  $(T,T,T,T,H,T,T,T,H,H,\dots)$.
  \item A macroeconomic series might record four numbers in each period,
  \[x_t=\big(\text{output}(t),\ \text{employment}(t),\ \text{debt}(t),\ \text{inflation}(t)\big),\]
  where $t=1,2,3,\dots$ denotes the time period.
\end{enumerate}
\src{L2, slide 22}
\end{example}

\begin{sourcenote}[Codomain of a sequence]
The slide's formal definition reads ``a function $x\colon\N\to\R$'', but its own examples have values that are not real
numbers (the letters H and T; four-number vectors). The general definition allows any codomain; sequences whose values
are real numbers are singled out next as \emph{numerical} sequences.
\end{sourcenote}

\begin{definition}[Numerical sequence][def:numerical-sequence]
A sequence $x=(x_1,x_2,x_3,\dots)$ with each $x_k\in\R$ is called a \term{numerical sequence}. \src{L2, slide 23}
\end{definition}

\begin{example}[Numerical sequences][ex:num-seq]
\begin{enumerate}
  \item The sequence $a=(7,7,7,7,\dots)$ is \emph{constant}.
  \item The $n$-th element of $b=(3,\,3.1,\,3.14,\,3.141,\dots)$ gives the first $n$ digits of the decimal expansion of $\pi$.
  \item The sequence $c=(-0.5,\,0.11,\,0.3596,\,0.365881,\dots)$ is given by $c_n=f(b_n)$, where $f(x)=x^2-9.5$. For example
  $c_2=3.1^2-9.5=9.61-9.5=0.11$.
  \item The $n$-th element of $d=(0,\,0.69,\,1.09,\,1.38,\,1.60,\,1.79,\,1.94,\dots)$ gives $\ln(n)$ truncated to two decimal
  places.
\end{enumerate}
\src{L2, slide 23}
\end{example}

\paragraph{Notation} All of the following may denote a sequence:
$x$, \ $x_1,x_2,x_3,\dots$, \ $(x_1,x_2,x_3,\dots)$, \ $(x_n)$, \ $(x_n)_{n\in\N}$, \ $(x_n)_{n=1}^{\infty}$.
The index need not be $n$; letters such as $i,j,k,l,m$ are also used. \src{L2, slide 23}

\section{Behaviour at infinity}\label{sec:at-infinity}

What happens to $a_n$ when $n$ grows arbitrarily large? Consider $a_n=2^{-n}$:
\[
(a_n)=\Big(\tfrac12,\tfrac14,\tfrac18,\tfrac1{16},\tfrac1{32},\tfrac1{64},\tfrac1{128},\tfrac1{256},\tfrac1{512},\tfrac1{1024},\tfrac1{2048},\dots\Big).
\]
As $n$ grows, $a_n$ stays positive but becomes smaller and smaller, and closer and closer to zero. So $0$ is a good
description of what happens ``at infinity''. If instead the sequence were
\[
(b_n)=\Big(\tfrac12,\tfrac14,\tfrac18,\tfrac1{16},\tfrac1{32},\tfrac1{64},\tfrac1{128},\tfrac1{256},\tfrac1{256},\tfrac1{256},\tfrac1{256},\dots\Big),
\]
then the number $\tfrac1{256}$ would better describe its behaviour at infinity. \src{L2, slide 24}

\begin{core}[Why a precise definition is needed]
The behaviour ``at infinity'' cannot be determined from any finite number of initial elements: $(a_n)$ and $(b_n)$ agree
in their first eight terms but head to different places. A definition is needed that speaks about \emph{all sufficiently
late} terms at once.
\end{core}

\section{The definition of a limit}\label{sec:limit-def}

\begin{definition}[Limit of a sequence][def:limit]
A number $x\in\R$ is a \term{limit} of a sequence $(a_n)$ if for any real number $\varepsilon>0$ there exists a natural
number $N$ such that for all $n>N$, $|a_n-x|<\varepsilon$. In this case we say that the sequence $(a_n)$ has a
\term{finite limit}, or that it \term{converges}. In symbols:
\[
\forall\varepsilon>0,\ \exists N\in\N:\ \forall n>N,\ |a_n-x|<\varepsilon. \tag*{\srcin{L2, slide 25}}
\]
\end{definition}

\begin{sourcenote}[A missing subscript]
On the slide, the symbolic version of the definition is printed with ``$|a-x|<\varepsilon$''; the subscript $n$ is missing.
The correct condition is $|a_n-x|<\varepsilon$, as in the verbal definition above it.
\end{sourcenote}

\subsection{Reading the definition}
\begin{core}[The limit as a game]
Read the definition from left to right, as a challenge and a response.
\begin{enumerate}[label=\arabic*.]
  \item \textbf{$\forall\varepsilon>0$}: a challenger names a tolerance $\varepsilon$, as small as they like.
  \item \textbf{$\exists N\in\N$}: you must reply with a point $N$ in the sequence. Your $N$ may depend on $\varepsilon$
  (it comes after $\varepsilon$; \cref{ex:quant-order}).
  \item \textbf{$\forall n>N$}: every term after position $N$ must then pass the test:
  \item \textbf{$|a_n-x|<\varepsilon$}: $a_n$ is within distance $\varepsilon$ of $x$ (\cref{sec:abs-distance}), i.e.\
  $x-\varepsilon<a_n<x+\varepsilon$.
\end{enumerate}
$x$ is a limit if you can win this game whatever $\varepsilon$ the challenger chooses. Equivalently: however narrow a band
is drawn around $x$, the sequence eventually enters the band and never leaves it.
\end{core}

\begin{figure}[ht]
\centering
\begin{tikzpicture}
\begin{axis}[stat, width=12cm, height=6cm, xmin=0, xmax=21, ymin=0.85, ymax=2.1, xlabel={$n$}, ylabel={$b_n$},
  xtick={1,5,10,15,20}, ytick={1,1.2,1.5,2}, yticklabels={$1$,$1+\varepsilon$,$1.5$,$2$}]
  \fill[accent!15] (axis cs:0,0.8) rectangle (axis cs:21,1.2);
  \addplot[guide, domain=0:21] {1.2};
  \addplot[guide, domain=0:21] {0.8};
  \addplot[thin, black, domain=0:21] {1};
  \addplot[only marks, mark=*, mark size=1.6pt, accent] coordinates {(1,2)(2,1.5)(3,1.3333)(4,1.25)};
  \addplot[only marks, mark=*, mark size=1.6pt, excol] coordinates {(5,1.2)(6,1.1667)(7,1.1429)(8,1.125)(9,1.1111)(10,1.1)(11,1.0909)(12,1.0833)(13,1.0769)(14,1.0714)(15,1.0667)(16,1.0625)(17,1.0588)(18,1.0556)(19,1.0526)(20,1.05)};
  \draw[corecol, thick] (axis cs:5,0.85) -- (axis cs:5,2.0) node[above, font=\small] {$N=5$};
\end{axis}
\end{tikzpicture}
\caption{The sequence $b_n=1+\frac1n$ with the band $|b_n-1|<\varepsilon$ for $\varepsilon=0.2$ (shaded). The terms $b_1,\dots,b_4$
lie outside the band; $b_5=1.2$ lies on its edge (since $|b_5-1|=0.2\not<0.2$); every term after $N=5$ lies strictly inside. A
smaller $\varepsilon$ requires a larger $N$, but for every $\varepsilon>0$ some $N$ works.}
\label{fig:epsilon-band}
\end{figure}

\section{Proving convergence}\label{sec:proving-convergence}

\begin{example}[Finding $N$ for particular $\varepsilon$][ex:2n-table]
Let $a_n=2^{-n}$ and $x=0$, so $|a_n-x|=|2^{-n}-0|=2^{-n}$. For each $\varepsilon$ below, find an $N$ such that
$|a_n-0|<\varepsilon$ for all $n>N$. \src{L2, slide 26}
\solution
\begin{center}
\begin{tabular}{@{}lll@{}}
\toprule
$\varepsilon$ & first term inside the band & a valid $N$\\
\midrule
$1$ & all terms ($a_1=\tfrac12<1$) & $N=1$\\
$0.1$ & $a_4=\tfrac1{16}=0.0625$ (since $a_3=0.125$) & $N=3$\\
$0.01$ & $a_7=\tfrac1{128}\approx0.0078$ (since $a_6\approx0.0156$) & $N=6$\\
$0.003$ & $a_9=\tfrac1{512}\approx0.00195$ (since $a_8\approx0.0039$) & $N=8$\\
\bottomrule
\end{tabular}
\end{center}
Because $2^{-n}$ decreases as $n$ increases, once one term is inside the band, all later terms are too.
\end{example}

\begin{note}[Any working $N$ will do]
$N=100$ works for $\varepsilon=1,\,0.1,\,0.01,\,0.003$ as well. It does not matter which particular $N$ we choose for which
$\varepsilon$: the important thing is that for \emph{any} $\varepsilon>0$ there is \emph{some} $N$ such that
$|a_n-x|<\varepsilon$ whenever $n>N$. In words, people sometimes say that $|a_n-x|<\varepsilon$ ``for large enough $n$'',
meaning the same statement. \src{L2, slide 27}
\end{note}

A table, however long, does not prove convergence: it deals with finitely many values of $\varepsilon$. A proof must handle an
\emph{arbitrary} $\varepsilon>0$.

\begin{example}[$2^{-n}$ converges to $0$][ex:2n-proof]
Prove that $x=0$ is a limit of $a_n=2^{-n}$. \src{L2, slide 28}
\solution
Consider an arbitrary $\varepsilon>0$. We want to show that for large enough $n$, $|2^{-n}-0|<\varepsilon$, that is, $2^{-n}<\varepsilon$.
Since $\log_2$ is increasing,
\[
2^{-n}<\varepsilon\ \Longleftrightarrow\ \log_2(2^{-n})<\log_2\varepsilon\ \Longleftrightarrow\ -n<\log_2\varepsilon\ \Longleftrightarrow\ n>-\log_2\varepsilon.
\]
Hence we can take any natural number $N>-\log_2\varepsilon$; then for all $n>N$,
\[
n>N>-\log_2\varepsilon\ \Longrightarrow\ |2^{-n}-0|<\varepsilon.
\]
Since $\varepsilon>0$ was arbitrary, $0$ is a limit of $(2^{-n})$. (There are many other proofs.) $\square$
\end{example}

\begin{external}[A proof without logarithms]
By Bernoulli's inequality (\cref{ex:bernoulli}), $2^n\ge n+1>n$, so $0<2^{-n}<\frac1n$. Given $\varepsilon>0$, choose a natural
number $N\ge\frac1\varepsilon$. For $n>N$: $2^{-n}<\frac1n<\frac1N\le\varepsilon$. This uses only the properties of the real
numbers and induction, and it shows a useful technique: bound the distance $|a_n-x|$ by something simpler (here $\frac1n$)
whose behaviour is easy to control.
\end{external}

\begin{example}[$1+\frac1n$ converges to $1$][ex:one-over-n]
Consider $b_n=1+\frac1n$, i.e.\ $(b_n)=\big(2,\tfrac32,\tfrac43,\tfrac54,\tfrac65,\tfrac76,\dots\big)$. Show that $x=1$ is a limit of
$(b_n)$. \src{L2, slide 29}
\solution
\begin{itemize}
  \item Consider an arbitrary $\varepsilon>0$. (We need to prove that the statement holds for \emph{any} $\varepsilon>0$.)
  \item Let $N=\lceil1/\varepsilon\rceil$ be the smallest natural number not less than $1/\varepsilon$. (It is enough to find \emph{one}
  $N$ that works for the given $\varepsilon$.)
  \item Consider an arbitrary $n>N$. (We again need to show that the inequality holds for \emph{any} $n>N$, given the $N$ we have
  picked.) Then
  \[
  |b_n-1|=\Big|\frac1n\Big|=\frac1n<\frac1N\le\varepsilon.
  \]
  \item Therefore $|b_n-1|<\varepsilon$ for all $n>N$, which we wanted to prove. $\square$
\end{itemize}
Each line of the proof corresponds to one quantifier of the definition: an arbitrary $\varepsilon$ ($\forall$), a chosen $N$
($\exists$), an arbitrary $n$ ($\forall$), and finally the inequality.
\end{example}

\begin{core}[Template for proving $a_n\to x$ from the definition]
\begin{enumerate}[label=\arabic*.]
  \item \emph{Scratch work:} simplify $|a_n-x|$ and find a condition on $n$ that makes it less than $\varepsilon$ (solve the
  inequality, or bound $|a_n-x|$ by something simpler first).
  \item \emph{Proof:} ``Let $\varepsilon>0$ be arbitrary. Choose $N=\dots$ (a natural number depending on $\varepsilon$). Let $n>N$.
  Then $|a_n-x|=\dots<\varepsilon$. Hence $x$ is a limit of $(a_n)$.''
\end{enumerate}
\end{core}

\begin{example}[A rational sequence][ex:rational-seq]
Prove from the definition that $a_n=\dfrac{n^2-1}{n^2+1}$ converges to $1$.
\solution
\emph{Scratch work.} $|a_n-1|=\Big|\frac{n^2-1-(n^2+1)}{n^2+1}\Big|=\frac{2}{n^2+1}<\frac2{n^2}\le\frac2n$ for $n\ge1$. So it is enough to
make $\frac2n<\varepsilon$, i.e.\ $n>\frac2\varepsilon$.

\emph{Proof.} Let $\varepsilon>0$ be arbitrary and choose a natural number $N\ge\frac2\varepsilon$. For $n>N$,
\[
|a_n-1|=\frac{2}{n^2+1}<\frac2n<\frac2N\le\varepsilon.
\]
Hence $a_n\to1$. $\square$
\end{example}

\section{Showing that a sequence has no limit}\label{sec:no-limit}

To show that a sequence does not converge, we must show that \emph{no} real number $x$ is a limit. Negating the definition
(\cref{sec:negation}): $x$ is \emph{not} a limit of $(a_n)$ when
\[
\exists\varepsilon>0:\ \forall N\in\N,\ \exists n>N:\ |a_n-x|\ge\varepsilon.
\]
Some tolerance $\varepsilon$ defeats every $N$: however late we start, some later term is at least $\varepsilon$ away from $x$.

\begin{example}[An oscillating sequence has no limit][ex:no-limit]
Consider the sequence $1,0,1,0,1,0,\dots$, that is, $c_n=1$ if $n$ is odd and $c_n=0$ if $n$ is even. Show that $(c_n)$ does
not have a (finite) limit. \src{L2, slide 30}
\solution
\begin{itemize}
  \item Towards a contradiction, assume that $x$ is a limit of $(c_n)$. (We should consider an \emph{arbitrary} $x$.)
  \item Consider $\varepsilon=\tfrac14$. (It is enough to find \emph{some} $\varepsilon>0$ that leads to a contradiction.)
  \item Then there should exist $N\in\N$ such that for all $n>N$, $|c_n-x|<\varepsilon$. (We should consider an arbitrary such $N$.)
  \item One of the numbers $c_{N+1},c_{N+2}$ is zero and the other is one; since $N+1,N+2>N$, it should be that
  $|1-x|<\tfrac14$ and $|0-x|<\tfrac14$. Then, by the triangle inequality (\cref{thm:triangle3}),
  \[
  1=|1-0|\le|1-x|+|x-0|<\tfrac14+\tfrac14=\tfrac12.
  \]
  \item We get $1<\tfrac12$, a contradiction. $\square$
\end{itemize}
\end{example}

The choice $\varepsilon=\tfrac14$ is not special: any $\varepsilon\le\tfrac12$ works. The idea is that the terms keep returning to
both $0$ and $1$, which are $1$ apart, so no single number can be within less than $\tfrac12$ of both; the triangle inequality
turns this picture into a proof.

\section{Further properties}\label{sec:limit-props}

\begin{external}[A sequence has at most one limit]
The slides say ``$x$ is \emph{a} limit''. In fact a sequence cannot have two different limits, so one may speak of
\emph{the} limit and write $\lim_{n\to\infty}a_n=x$ or $a_n\to x$. \emph{Proof.} Suppose $x\ne y$ are both limits and take
$\varepsilon=\frac{|x-y|}{2}>0$. There are $N_1$ and $N_2$ with $|a_n-x|<\varepsilon$ for $n>N_1$ and $|a_n-y|<\varepsilon$ for $n>N_2$. For
$n>\max\{N_1,N_2\}$ both hold, and by the triangle inequality $|x-y|\le|x-a_n|+|a_n-y|<2\varepsilon=|x-y|$, a contradiction.
\end{external}

The following results are in the style of the first question of the January 2025 paper (as summarised in the study guide).
They show how the definition is used to prove general facts.

\begin{example}[Limits preserve non-negativity][ex:limit-nonneg]
Let $a_n\to a$ and suppose $a_n\ge0$ for every $n$. Prove that $a\ge0$. Hint: assume, towards a contradiction, that $a<0$.
\src{SG: Jan 2025 Q1(b)}
\solution
Suppose, for a contradiction, that $a<0$. Take $\varepsilon=-a>0$ (the distance from $a$ to $0$). By the definition of the limit there
is $N$ with $|a_n-a|<-a$ for all $n>N$. Take $n=N+1$. Then $a_n-a\le|a_n-a|<-a$, so $a_n<0$. This contradicts $a_n\ge0$. Hence
$a\ge0$. (The negation of $a\ge0$ is $a<0$ by trichotomy, \cref{sec:negation}.)
\end{example}

\begin{example}[Comparing limits][ex:limit-compare]
Let $b_n\to b$ and $c_n\to c$ with $c_n\ge b_n$ for every $n$. Prove that $c\ge b$. \src{SG: Jan 2025 Q1(c)}
\solution
Let $a_n=c_n-b_n\ge0$. First we show $a_n\to c-b$. Given $\varepsilon>0$, choose $N_1$ with $|b_n-b|<\varepsilon/2$ for $n>N_1$, and $N_2$
with $|c_n-c|<\varepsilon/2$ for $n>N_2$. For $n>N=\max\{N_1,N_2\}$, by the triangle inequality,
\[
|a_n-(c-b)|=|(c_n-c)-(b_n-b)|\le|c_n-c|+|b_n-b|<\tfrac\varepsilon2+\tfrac\varepsilon2=\varepsilon.
\]
So $a_n\to c-b$. Since $a_n\ge0$ for all $n$, \cref{ex:limit-nonneg} gives $c-b\ge0$, i.e.\ $c\ge b$.
\end{example}

\begin{warning}[Strict inequalities need not survive the limit]
If $a_n>0$ for every $n$ and $a_n\to a$, it does \emph{not} follow that $a>0$: the sequence $a_n=\frac1n$ is always positive
but its limit is $0$. Only the non-strict inequality $a\ge0$ is guaranteed.
\end{warning}

\section{Limits in economics}\label{sec:limits-econ}

\begin{example}[Convergence to a steady state][ex:steady-state]
A simple model of capital accumulation sets $K_{t+1}=0.8K_t+20$, with $K_1=50$. Show that $K_t=100-50\cdot(0.8)^{t-1}$ and that
$K_t\to100$.
\solution
The formula holds for $t=1$ ($100-50=50$). If it holds for $t$, then
$K_{t+1}=0.8\big(100-50(0.8)^{t-1}\big)+20=100-50(0.8)^t$, which is the formula for $t+1$; so by induction it holds for all $t$.
Then $|K_t-100|=50(0.8)^{t-1}$. Given $\varepsilon>0$, this is less than $\varepsilon$ once $(0.8)^{t-1}<\frac\varepsilon{50}$, i.e.\
$t-1>\frac{\ln(\varepsilon/50)}{\ln0.8}$ (dividing by $\ln0.8<0$ reverses the inequality). Choosing $N$ larger than
$1+\frac{\ln(\varepsilon/50)}{\ln0.8}$ completes the proof. The capital stock approaches the \emph{steady state} $K^*=100$, the value
at which $K^*=0.8K^*+20$.
\end{example}

\begin{external}[Infinite sums are limits]
The present value of a perpetuity paying $\pounds1$ per year forever at interest rate $i>0$ is defined as the limit of the
present values of finite streams, $\sum_{t=1}^{n}(1+i)^{-t}=\frac{1-(1+i)^{-n}}{i}$ (\cref{ex:geometric}). Since $(1+i)^{-n}\to0$, the
limit is $\frac1i$. Infinite sums of this kind are the subject of Lecture~3.
\end{external}

% --------------------------------------------------------------------
\section{Chapter review}
% --------------------------------------------------------------------
\subsection*{Summary}
A sequence is an infinite ordered list indexed by $\N$, formally a function with domain $\N$; a numerical sequence has real
values. Its behaviour at infinity cannot be read from finitely many terms, which is why limits are defined with quantifiers:
$x$ is a limit of $(a_n)$ if for every $\varepsilon>0$ there is $N\in\N$ such that $|a_n-x|<\varepsilon$ for all $n>N$. To prove
convergence, take an arbitrary $\varepsilon$, produce an $N$ (which may depend on $\varepsilon$), take an arbitrary $n>N$ and verify the
inequality. To prove that there is no limit, assume a limit $x$ exists and use one well-chosen $\varepsilon$ and the triangle
inequality to reach a contradiction. Limits are unique, preserve non-strict inequalities, and describe long-run behaviour such
as convergence to a steady state.

\subsection*{Key definitions}
\begin{itemize}
  \item \textbf{Sequence:} $(x_1,x_2,\dots)$, formally a function on $\N$ with $x_n=x(n)$; \textbf{numerical} if $x_n\in\R$.
  \item \textbf{Limit:} $x$ is a limit of $(a_n)$ if $\forall\varepsilon>0\ \exists N\in\N\ \forall n>N:\ |a_n-x|<\varepsilon$; the
  sequence then \textbf{converges} (has a finite limit).
\end{itemize}

\subsection*{Key statements}
\begin{gather*}
\text{$x$ is a limit:}\quad \forall\varepsilon>0,\ \exists N\in\N:\ \forall n>N,\ |a_n-x|<\varepsilon;\\
\text{$x$ is not a limit:}\quad \exists\varepsilon>0:\ \forall N\in\N,\ \exists n>N:\ |a_n-x|\ge\varepsilon;\\
2^{-n}\to0,\qquad 1+\tfrac1n\to1,\qquad (1,0,1,0,\dots)\ \text{has no limit.}
\end{gather*}

\subsection*{Connections}
The definition combines nearly everything in the book: quantifiers and their order (\cref{ch:logic}), sets and functions
(\cref{ch:sets,ch:func}), the properties of $\R$ (\cref{ch:numbers}), solving inequalities for $n$ (\cref{ch:ineq}), distance and the
triangle inequality (\cref{ch:abs}) and induction (\cref{ch:ind}). It is the foundation of Topic~3: properties and calculation of
limits, infinite sums, continuity and derivatives are all defined through limits.

\subsection*{Common mistakes}
\begin{itemize}
  \item Choosing $N$ before $\varepsilon$, or a single $N$ meant to work for every $\varepsilon$.
  \item ``Proving'' convergence with a table of values for a few $\varepsilon$.
  \item Fixing a particular $n$ instead of an arbitrary $n>N$.
  \item Writing $|a_n-x|\le\varepsilon$ in place of $<\varepsilon$ without care (harmless in practice, but match the definition).
  \item In a non-existence proof, choosing a particular $x$ instead of an arbitrary one.
  \item Concluding $a>0$ from $a_n>0$.
\end{itemize}

\subsection*{Exam checklist}
\begin{checklist}
  \item Define a sequence and a numerical sequence; give examples.
  \item State the definition of a limit exactly, with quantifiers in the right order.
  \item For a given sequence and $\varepsilon$, find a valid $N$.
  \item Prove from the definition that a given sequence converges to a given number.
  \item Prove that a given sequence has no limit.
  \item Prove that limits preserve $\ge$, using contradiction.
\end{checklist}

% --------------------------------------------------------------------
\section{Practice questions}
% --------------------------------------------------------------------
\pq[basic] Write down the first five terms of $a_n=\frac{(-1)^n}{n}$, $b_n=\frac{n}{n+1}$ and $c_n=n^2-3n$.

\pq[basic] For $a_n=\frac1n$ and $x=0$, find an $N$ that works for $\varepsilon=0.1$, for $\varepsilon=0.01$ and for $\varepsilon=0.003$.

\pq[conceptual] Explain why the following is \emph{not} a valid definition of $a_n\to x$: ``there is $N\in\N$ such that for every
$\varepsilon>0$ and every $n>N$, $|a_n-x|<\varepsilon$.'' Which sequences satisfy it?

\pq[mathematical] Prove from the definition that $\frac{3n+1}{n}\to3$.

\pq[mathematical] Prove from the definition that $\frac{(-1)^n}{n}\to0$.

\pq[mathematical] Prove that the sequence $a_n=(-1)^n$ has no limit.

\pq[application] A savings balance evolves as $B_{t+1}=0.5B_t+10$ with $B_1=30$. Show by induction that $B_t=20+10\cdot(0.5)^{t-1}$,
and prove from the definition that $B_t\to20$.

\pq[multi-step] Let $a_n\to a$. Prove that $3a_n\to3a$, directly from the definition.

\pq[exam-style] (a) Define what it means that a sequence $(a_n)$ has a finite limit $a$.
(b) Let $a_n\to a$, and suppose $a_n\le5$ for all $n$. Prove that $a\le5$. Hint: assume, towards a contradiction, that $a>5$.
(c) Give an example of a sequence with $a_n<5$ for all $n$ whose limit is not less than $5$.

% --------------------------------------------------------------------
\section{Solutions to practice questions}
% --------------------------------------------------------------------
\pqsol{10.1}
$a$: $-1,\tfrac12,-\tfrac13,\tfrac14,-\tfrac15$. \ $b$: $\tfrac12,\tfrac23,\tfrac34,\tfrac45,\tfrac56$. \ $c$: $-2,-2,0,4,10$.

\pqsol{10.2}
We need $\frac1n<\varepsilon$ for $n>N$, i.e.\ $n>\frac1\varepsilon$. So $N=10$ works for $\varepsilon=0.1$ (for $n\ge11$, $\frac1n<0.1$); $N=100$ for
$\varepsilon=0.01$; and $N=333$ for $\varepsilon=0.003$ (since $\frac1{334}<0.003<\frac1{333}$). Any larger $N$ also works.

\pqsol{10.3}
Here $N$ is chosen \emph{before} $\varepsilon$ and must work for every $\varepsilon$ simultaneously. If $|a_n-x|<\varepsilon$ for every
$\varepsilon>0$, then $|a_n-x|$ is a non-negative number smaller than every positive number, so it is $0$. The condition therefore says
that $a_n=x$ for all $n>N$: only sequences that are eventually constant equal to $x$ satisfy it. The correct definition lets $N$
depend on $\varepsilon$ (\cref{ex:quant-order}).

\pqsol{10.4}
$\big|\frac{3n+1}n-3\big|=\frac1n$. Given $\varepsilon>0$, choose a natural number $N\ge\frac1\varepsilon$. For $n>N$:
$\frac1n<\frac1N\le\varepsilon$. Hence the limit is $3$.

\pqsol{10.5}
$\big|\frac{(-1)^n}{n}-0\big|=\frac1n$. Given $\varepsilon>0$, choose $N\ge\frac1\varepsilon$; for $n>N$, $\frac1n<\varepsilon$.

\pqsol{10.6}
Suppose, for a contradiction, that $x$ is a limit, and take $\varepsilon=\tfrac12$. There is $N$ with $|a_n-x|<\tfrac12$ for $n>N$. Among
$N+1$ and $N+2$ one is even and one is odd, so $|1-x|<\tfrac12$ and $|-1-x|<\tfrac12$. By the triangle inequality,
$2=|1-(-1)|\le|1-x|+|x-(-1)|<\tfrac12+\tfrac12=1$, a contradiction. So $(a_n)$ has no limit.

\pqsol{10.7}
Base: $20+10=30=B_1$. Step: if $B_t=20+10(0.5)^{t-1}$ then $B_{t+1}=0.5\big(20+10(0.5)^{t-1}\big)+10=20+10(0.5)^t$.
Limit: $|B_t-20|=10(0.5)^{t-1}=20\cdot2^{-t}$. By Bernoulli, $2^t>t$, so $|B_t-20|<\frac{20}{t}$. Given $\varepsilon>0$, choose $N\ge\frac{20}\varepsilon$;
for $t>N$, $|B_t-20|<\frac{20}t<\frac{20}N\le\varepsilon$.

\pqsol{10.8}
Let $\varepsilon>0$. Since $a_n\to a$, there is $N$ with $|a_n-a|<\frac\varepsilon3$ for all $n>N$. For such $n$,
$|3a_n-3a|=3|a_n-a|<3\cdot\frac\varepsilon3=\varepsilon$. Hence $3a_n\to3a$. (The trick is to apply the definition for $a_n$ with the smaller
tolerance $\varepsilon/3$.)

\pqsol{10.9}
(a) $(a_n)$ has finite limit $a\in\R$ if for every $\varepsilon>0$ there exists $N\in\N$ such that $|a_n-a|<\varepsilon$ for all $n>N$.
(b) Suppose $a>5$ and let $\varepsilon=a-5>0$. There is $N$ with $|a_n-a|<a-5$ for all $n>N$. For $n=N+1$:
$a-a_n\le|a_n-a|<a-5$, so $a_n>5$, contradicting $a_n\le5$. Hence $a\le5$.
(c) $a_n=5-\frac1n$: every term is less than $5$, but $a_n\to5$ (since $|a_n-5|=\frac1n$), and $5$ is not less than $5$.
